% SPDX-License-Identifier: Apache-2.0
\section{A General I2C Master}
\label{sec:x-i2c}

I2C is two wires and a protocol. Both are open drain: every device may
pull a line low and nobody may drive it high, so a line nobody pulls
is high through the board's pull-up. What that buys is a bus where any
device can hold the clock, and a master that has to watch the lines
rather than count cycles at them.

\code{I2cInit}, which \code{//docs:hdmi} describes, replays a fixed
list of writes into one chip and cannot read. \code{I2c} is the other
kind: four registers, and one command at a time.

\begin{center}\footnotesize
\begin{tabular}{@{}llp{0.52\columnwidth}@{}}
\toprule
Offset & Name & What it is \\
\midrule
\code{0x00} & \code{ctrl} & the divider, and the interrupt enable \\
\code{0x04} & \code{cmd} & one piece of a transaction; writing it
starts the piece \\
\code{0x08} & \code{data} & the byte a read took \\
\code{0x0c} & \code{state} & busy, done, not acknowledged, arbitration
lost \\
\bottomrule
\end{tabular}
\end{center}

A command is at most a start, one byte either way, and a stop. That is
not a limitation of the hardware so much as where the line between
hardware and software falls: the shape of a transaction is a
program's business, and every I2C transaction is a sequence of these
pieces. Reading a register of a device is four commands: address it
for writing, write the register number, address it again for reading
after a repeated start, and read the byte with a NACK and a stop,
because a master says NACK to the last byte it wants.

\subsection{What the Master Watches}

Two of the four states in \code{state} are there because the bus is
shared, and both are read from the lines rather than assumed.

A device that is not ready holds the clock low after the master
releases it. That is clock stretching, and the master waits: the
quarter of a bit does not advance while the clock reads low against
what the master drives. A master that counted cycles instead would
walk over the device.

A data line that reads low where the master released it is another
master driving the bus, which is arbitration lost. The command ends
there and the bit is set, because two masters that both keep going
corrupt each other's bytes.

There is a third thing the lines must be watched for, and it is the
one this example was written twice for: between the pieces of one
transaction, the clock is parked \emph{low} and not high. A data line
that moves while the clock is high is a start or a stop to every
device on the bus, so a master that parks the clock high and then
loads the next byte has sent a start nobody asked for. The device
model catches it: the byte after such a park arrives rotated by one
bit, because the device took the spurious start as the beginning of a
new address.

\subsection{What Checks It}

\code{I2cDev}, a model of a device with a byte per register, stepped
between cycles from the loop as the I2C device of \code{//docs:hdmi}
is. It acknowledges its own address and nothing else, answers reads
from the register its pointer names, and can be told to hold the clock
low for a few cycles after every byte.

Five tests in \code{//lib/parts:parts\_test} drive the master against
it: a write of two bytes under one address, a read after a repeated
start, an address nobody answers, which sets the NACK bit and clears
it again, and a device that stretches the clock, which changes only
how long the transaction takes.

\begin{figure}[htbp]
\centering
\input{i2c_timing.tex}
\caption{The address going out: the clock and the data as the master
pulls them, the lines as they read back, the step through the
transaction, and the byte shifting out a bit per clock.}
\label{fig:x-i2c}
\end{figure}

\lstinputlisting[caption={\code{ex\_i2c.rs}. Run with \code{bazel run //lib/examples:ex\_i2c}.},label={lst:x-i2c}]{lib/examples/ex_i2c.rs}

\lstinputlisting[style=ascii,caption={What \code{ex\_i2c} printed when the build ran it: a register read out of the device, and an address nobody answered.},label={lst:x-i2c-out}]{out_i2c.txt}
